\documentclass[10pt,a4paper]{amsart}
\usepackage{amsmath,amssymb,amsthm,graphicx}
% Retained preamble for the Stacks Project chapter {CH} (shared project preamble).
% The chapter file pinned-original-source.tex only carries \input{{preamble}} on its first line,
% so no preamble text lives inside the chapter itself.  The Stacks Project uses one single shared
% preamble for every chapter: the snapshot stages it three times (sources/stacks-chow/preamble.tex,
% sources/stacks-etale-cohomology/preamble.tex and focused-seeds/AGP/sources/stacks-snapshot-20261011/
% preamble.tex) and those copies are identical in size.  The block below is copied mechanically,
% line for line, from sources/stacks-chow/preamble.tex lines 14-18 and 154-217: the Xy-pic setup and
% the native theorem declarations together with the mathematical macros.  Nothing was added, removed
% or redefined; the wrapper class and the amsmath/amssymb/amsthm/graphicx base packages are supplied
% by the assembler.
% For commutative diagrams we use Xy-pic
% Chapter: stacks-spaces-more-morphisms.tex -- More on Morphisms of Spaces.
\usepackage[all]{xy}

% We use 2cell for 2-commutative diagrams.
\xyoption{2cell}

% Theorem environments.
%
\theoremstyle{plain}
\newtheorem{theorem}[subsection]{Theorem}
\newtheorem{proposition}[subsection]{Proposition}
\newtheorem{lemma}[subsection]{Lemma}

\theoremstyle{definition}
\newtheorem{definition}[subsection]{Definition}
\newtheorem{example}[subsection]{Example}
\newtheorem{exercise}[subsection]{Exercise}
\newtheorem{situation}[subsection]{Situation}

\theoremstyle{remark}
\newtheorem{remark}[subsection]{Remark}
\newtheorem{remarks}[subsection]{Remarks}

\numberwithin{equation}{subsection}

% Macros
%
\def\lim{\mathop{\mathrm{lim}}\nolimits}
\def\colim{\mathop{\mathrm{colim}}\nolimits}
\def\Spec{\mathop{\mathrm{Spec}}}
\def\Hom{\mathop{\mathrm{Hom}}\nolimits}
\def\Ext{\mathop{\mathrm{Ext}}\nolimits}
\def\SheafHom{\mathop{\mathcal{H}\!\mathit{om}}\nolimits}
\def\SheafExt{\mathop{\mathcal{E}\!\mathit{xt}}\nolimits}
\def\Sch{\mathit{Sch}}
\def\Mor{\mathop{\mathrm{Mor}}\nolimits}
\def\Ob{\mathop{\mathrm{Ob}}\nolimits}
\def\Sh{\mathop{\mathit{Sh}}\nolimits}
\def\NL{\mathop{N\!L}\nolimits}
\def\CH{\mathop{\mathrm{CH}}\nolimits}
\def\proetale{{pro\text{-}\acute{e}tale}}
\def\etale{{\acute{e}tale}}
\def\QCoh{\mathit{QCoh}}
\def\Ker{\mathop{\mathrm{Ker}}}
\def\Im{\mathop{\mathrm{Im}}}
\def\Coker{\mathop{\mathrm{Coker}}}
\def\Coim{\mathop{\mathrm{Coim}}}

% Boxtimes
%
\DeclareMathSymbol{\boxtimes}{\mathbin}{AMSa}{"02}

%
% Macros for moduli stacks/spaces
%
\def\QCohstack{\mathcal{QC}\!\mathit{oh}}
\def\Cohstack{\mathcal{C}\!\mathit{oh}}
\def\Spacesstack{\mathcal{S}\!\mathit{paces}}
\def\Quotfunctor{\mathrm{Quot}}
\def\Hilbfunctor{\mathrm{Hilb}}
\def\Curvesstack{\mathcal{C}\!\mathit{urves}}
\def\Polarizedstack{\mathcal{P}\!\mathit{olarized}}
\def\Complexesstack{\mathcal{C}\!\mathit{omplexes}}
% \Pic is the operator that assigns to X its picard group, usage \Pic(X)
% \Picardstack_{X/B} denotes the Picard stack of X over B
% \Picardfunctor_{X/B} denotes the Picard functor of X over B
\def\Pic{\mathop{\mathrm{Pic}}\nolimits}
\def\Picardstack{\mathcal{P}\!\mathit{ic}}
\def\Picardfunctor{\mathrm{Pic}}
\def\Deformationcategory{\mathcal{D}\!\mathit{ef}}
\usepackage{cleveref}
\crefname{originallemma}{original lemma}{original lemmas}
\Crefname{originallemma}{Original lemma}{Original lemmas}
\makeatletter
\newlabel{lemma-where-closed-immersion}{{05XA}{source}{}{sourcepointer-1}{}}
\newlabel{lemma-where-closed-immersion@cref}{{[originallemma][0][]05XA}{[1][0][]source}{}{}{}}
\makeatother
\crefname{originallemma}{original lemma}{original lemmas}
\Crefname{originallemma}{Original lemma}{Original lemmas}
\makeatletter
\newlabel{lemma-where-surjective-flat}{{05XC}{source}{}{sourcepointer-2}{}}
\newlabel{lemma-where-surjective-flat@cref}{{[originallemma][0][]05XC}{[1][0][]source}{}{}{}}
\makeatother
\crefname{originalreference}{original reference}{original references}
\Crefname{originalreference}{Original reference}{Original references}
\makeatletter
\newlabel{morphisms-lemma-characterize-flat-closed-immersions}{{04PW}{source}{}{sourcepointer-3}{}}
\newlabel{morphisms-lemma-characterize-flat-closed-immersions@cref}{{[originalreference][0][]04PW}{[1][0][]source}{}{}{}}
\makeatother
\title{The locus where a proper flat morphism is an isomorphism is open}
\author{The Stacks Project Authors}
\date{stacks-project:spaces-more-morphisms.tex:local-snapshot-20261011; GNU FDL 1.2 or later}
\begin{document}
\maketitle
\noindent\textit{Editorial note.} Counted span is the author's own proof environment, source lines 11814-11833, immediately following the statement at 11791-11812.\par\medskip
\noindent\textit{Editorial note.} The author's parenthetical remark at the end of the proof is retained verbatim.\par\medskip
\noindent\textit{Editorial note.} Two same-chapter lemmas and one cross-chapter lemma are recorded in assembly.external\_references.\par\medskip
\noindent\textit{Allowed standard prerequisites (editorial).}\begin{itemize}
\item The locus where a morphism is a closed immersion. chapter lemma lemma-where-closed-immersion, cited at proof line 11820. Supplies an open subspace W\_2 of Z factoring exactly the base changes for which f becomes a closed immersion.
\item The locus where a morphism is surjective and flat. chapter lemma lemma-where-surjective-flat, cited in the proof. Supplies an open subspace W\_1 factoring exactly the base changes for which f becomes surjective flat.
\item A surjective flat closed immersion of schemes is an isomorphism. Morphisms, lemma morphisms-lemma-characterize-flat-closed-immersions, cited at proof lines 11829-11830. Converts 'surjective flat closed immersion' into 'isomorphism' for the representing scheme morphism.
\end{itemize}\medskip
\ifdefined\hypertarget\hypertarget{sourcepointer-1}{}\fi
\noindent\textit{Source reference.} Original lemma 05XA: sources/stacks-spaces-more-morphisms/pinned-original-source.tex, source lines 11642--11642. Navigation pointer only: the counted proof cites this statement by \textbackslash{}ref as a standing input; the deduction the counted proof has to make is carried out inside the counted span, and no part of that deduction is delegated to this pointer.\par\medskip
\ifdefined\hypertarget\hypertarget{sourcepointer-2}{}\fi
\noindent\textit{Source reference.} Original lemma 05XC: sources/stacks-spaces-more-morphisms/pinned-original-source.tex, source lines 11729--11729. Navigation pointer only: the counted proof cites this statement by \textbackslash{}ref as a standing input; the deduction the counted proof has to make is carried out inside the counted span, and no part of that deduction is delegated to this pointer.\par\medskip
\ifdefined\hypertarget\hypertarget{sourcepointer-3}{}\fi
\noindent\textit{Source reference.} Original reference 04PW: focused-seeds/geometry/S9-spaces-where-isomorphism-spaces/ref-tags.txt, source lines 4--4. Navigation pointer only: the counted proof cites this statement by \textbackslash{}ref as a standing input; the deduction the counted proof has to make is carried out inside the counted span, and no part of that deduction is delegated to this pointer.\par\medskip
\begin{lemma}
\label{lemma-where-isomorphism}
Consider a commutative diagram
$$
\xymatrix{
X \ar[rr]_f \ar[rd]_p & & Y \ar[ld]^q \\
& Z
}
$$
of algebraic spaces. Assume that
\begin{enumerate}
\item $p$ is locally of finite presentation,
\item $p$ is flat,
\item $p$ is universally closed,
\item $q$ is locally of finite type,
\item $q$ is closed, and
\item $q$ is separated.
\end{enumerate}
Then there exists an open subspace $W \subset Z$
such that a morphism $Z' \to Z$ factors through $W$ if and only if the
base change $f_{Z'} : X_{Z'} \to Y_{Z'}$ is an isomorphism.
\end{lemma}

\begin{proof}
By
Lemma \ref{lemma-where-surjective-flat}
there exists an open subspace $W_1 \subset Z$ such that
$f_{Z'}$ is surjective and flat if and only if $Z' \to Z$
factors through $W_1$. By
Lemma \ref{lemma-where-closed-immersion}
there exists an open subspace $W_2 \subset Z$ such that
$f_{Z'}$ is a closed immersion if and only if $Z' \to Z$
factors through $W_2$. We claim that $W = W_1 \cap W_2$ works.
Certainly, if $f_{Z'}$ is an isomorphism, then $Z' \to Z$
factors through $W$. Hence it suffices to show that
$f_W$ is an isomorphism. By construction $f_W$ is a
surjective flat closed immersion. In particular $f_W$ is
representable. Since a surjective flat closed immersion of
schemes is an isomorphism (see
Morphisms, Lemma \ref{morphisms-lemma-characterize-flat-closed-immersions})
we win. (Note that actually $f_W$ is locally of finite presentation,
whence open, so you can avoid the use of this lemma if you like.)
\end{proof}

\end{document}
